

\section{Weighted Activity Selection}\label{dp: sec: activity selection}\doHeader{Activity Selection}

\subsection*{Warm-up --- counting pairwise-disjoint subsets}

\begin{mylist}
\item[\bf input:]
  A collection $C=\{[s_1, e_1], [s_2, e_2], \ldots, [s_n, e_n]\}$ of intervals.

  For convenience we assume $e_1 \le e_2 \le \cdots \le e_n$.  (If not, sort them.)

\item[\bf output:]
  The number of subsets $C'\subseteq C$ of $C$ that are pairwise disjoint.

  (These include the subsets of size one or zero.)
\end{mylist}

Here is an example input, $C$, consisting of seven intervals.  (The interval heights are irrelevant.)

\begin{center}
  \begin{tikzpicture}[scale=0.5]\footnotesize
    \draw (5,0) -- ++(-2,0) node[anchor=south] {1} -- ++(-2,0); 
    \draw (7,1) -- ++(-2,0) node[anchor=south] {2} -- ++(-2,0); 
    \draw (8,0) -- ++(-1,0) node[anchor=south] {3} -- ++(-1,0); 
    \draw (10,3) -- ++(-2.5,0) node[anchor=south] {4} -- ++(-2.5,0); 
    \draw (12,2) -- ++(-2.75,0) node[anchor=south] {5} -- ++(-2.75,0); 
    \draw (13,0) -- ++(-2,0) node[anchor=south] {6} -- ++(-2,0); 
    \draw (14,1) -- ++(-1.5,0) node[anchor=south] {7} -- ++(-1.5,0); 
  \end{tikzpicture}
\end{center}

Following \prob{Activity Selection},
call each interval in $C$ a \emph{job}, and interpret each job $[s_i, e_i]$ as
an activity that lasts from start time $s_i$ to end time $e_i$.  A pairwise-disjoint subset of $C$
is then just a collection of jobs can be done without doing more than one job at any time.

Our first goal is to come up with some recursive or inductive understanding
of the collection of pairwise-disjoint subsets.  To start, partition the pairwise-disjoint subsets into two types:
\begin{mylist}
\item[(A)] Pairwise-disjoint subsets of $C$ that don't contain the last job $[s_n, e_n]$. 
\item[(B)] Pairwise-disjoint subsets of $C$ that do contain the last job $[s_n, e_n]$. 
\end{mylist}

The subsets of type (A) are just the pairwise-disjoint
subsets of the first $n-1$ jobs,
that is, of the collection $C_A=\{[s_1, e_1], [s_2, e_2], \ldots, [s_{n-1}, e_{n-1}]\}$.


What about the subsets of type (B)?
Any subset of type (B) contains the job $[s_n, e_n]$,
so can't contain any job that overlaps with job $[s_n, e_n]$,
that is, can't contain any job $[s_i, e_i]$ with $i \ge i'$,
where $i' = \min\{ i: e_i \ge s_n\}$
is the index of the first job that overlaps job $[s_n, e_n]$.

In the example, $n=7$ and $i'=5$.
The jobs overlapping the last job $[s_n, e_n]$ are shown in blue:

\begin{center}
\begin{tikzpicture}[scale=0.5]\footnotesize
   \draw (5,0) -- ++(-2,0) node[anchor=south] {1} -- ++(-2,0); 
   \draw (7,1) -- ++(-2,0) node[anchor=south] {2} -- ++(-2,0); 
   \draw (8,0) -- ++(-1,0) node[anchor=south] {3} -- ++(-1,0); 
   \draw (10,3) -- ++(-2.5,0) node[anchor=south] {4} -- ++(-2.5,0); 
   \draw[color=blue,dashed] (12,2) -- ++(-2.75,0) node[anchor=south] {$i'=5$} -- ++(-2.75,0); 
   \draw[color=blue,dashed] (13,0) -- ++(-2,0) node[anchor=south] {6} -- ++(-2,0); 
   \draw[color=blue] (14,1) -- ++(-1.5,0) node[anchor=south] {$n=7$} -- ++(-1.5,0); 
\end{tikzpicture}
\end{center}

So, the subsets of type (B) consist of $[s_n, e_n]$,
together with any pairwise-disjoint subset of the first $i'-1$ jobs,
that is, of the collection $C_B=\{[s_1, e_1], [s_2, e_2], \ldots, [s_{i'-1}, e_{i'-1}]\}$.
The collection $C_B$ can be empty.
(In the example, $C_B$ contains the first four jobs.)

This approach gives the following dynamic-programming algorithm.

\paragraph{subproblems and recurrence.}
For $j\in\{0,1,2,\ldots, n\}$,
define $N(j)$ to be the number of pairwise-disjoint subsets of 
$\{[s_1, e_1], [s_2, e_2], \ldots, [s_j, e_j]\}$.
The output will be $N(n)$.
The recurrence relation is
\[
  N(j) = \begin{cases}
    1 & \text{ if } j=0 \\
    N(j-1) + N(i'-1), \text{ where }  i' = \min\{ i : e_i \ge s_j \}
    & \text{ if } j>0.
  \end{cases} 
\]

\paragraph{correctness.}
The recurrence relation follows from the preceding discussion.
We don't give a careful proof here.

\paragraph{run time.}
% [review 2026-09-27] fix: N(j) is defined for j in {0,...,n}, so there are n+1 subproblems, not n
There are $n+1$ subproblems.  The time to solve each depends on how we compute $i'$.
Here is the iterative version of the algorithm:

\begin{myframed}
  \begin{steps}
  \item[]
    \textsf{countSchedules}$(C=\{[s_1, e_1], [s_2, e_2], \ldots, [s_n, e_n]\})$:
    \comment{--- iterative version}
    \step for $j=0,1,\ldots, n$:
    \step~~~~set \(M[j]  = \begin{cases}
      1 & \text{ if } j=0 \\
      M[j-1] + M[i'-1], \text{ where }  i' = \min\{ i : e_i \ge s_j \}
      & \text{ if } j>0
    \end{cases}
    \).
    \step return $M[n]$
  \end{steps}
\end{myframed}

If we assume constant-time arithmetic (which is not realistic here!)
the algorithm can be implemented to run in $\Theta(n\log n)$ time.
The main obstacle is computing $i'$ in each iteration.
% [review 2026-09-27] fix: i' = min{i : e_i >= s_j} is found by binary search over the end times (already sorted), not over jobs sorted by start times
Since the jobs are sorted by end times,
$i'$ can be computed in each iteration in $O(\log n)$ time using binary search.
Then the $n$ iterations take total time $\Theta(n\log n)$, as does the sort.

But the value of $M[i]$ can be as large as $2^i$,
so in the worst case requires $n$ bits to store.
Then each addition takes $\Theta(n)$ time,
so each of the $n$ iterations takes $\Theta(n)$ time,
for a total time of $\Theta(n^2)$.
Here is what we have figured out:
\begin{theorem}
  Given a collection $C$ of $n$ intervals,
  the number of pairwise-disjoint subsets of $C$
  can be computed in $O(n^2)$ time.
\end{theorem}

\subsection*{Weighted Activity Selection}

\begin{mylist}
\item[\bf input:]
  A collection $C=\{[s_1, e_1], [s_2, e_2], \ldots, [s_n, e_n]\}$ of intervals (called \emph{jobs}),
  and, for each job $[s_i, e_i]$, an associated value $v_i \ge 0$.
  Again assume $e_1 \le e_2 \le \cdots \le e_n$.

\item[\bf output:]
  The maximum, over all pairwise-disjoint subsets $C'$ of $C$,
  of the total value of intervals in $C'$, that is $\sum_{[s_i, e_i]\in C'} v_i$.
\end{mylist}

\paragraph{subproblems and recurrence relation.} 
Fix an instance.
For $j\in\{0,1,\ldots,n\}$,
define $V(j)$ to be the maximum value of any pairwise-disjoint subset from
the first $j$ intervals
$\{[s_1, e_1], [s_2, e_2], \ldots, [s_j, e_j]\}$.
The final answer is $V(n)$.
The recurrence relation is
\[
  V(j) = \begin{cases}
    0 & \text{ if } j=0 \\
    \max\big( V(j-1), V(i'-1) + v_j\big), \text{ where }  i' = \min\{ i : e_i \ge s_j \},
    & \text{ if } j>0.
  \end{cases} 
\]

\paragraph{correctness.}
As before, we divide the pairwise-disjoint subsets of $\{[s_1, e_1], [s_2, e_2], \ldots, [s_j, e_j]\}$ into two types:
\begin{mylist}
\item[(A)] The pairwise-disjoint subsets that don't include the last job $[s_j, e_j]$. 
\item[(B)] The pairwise-disjoint subsets that do include the last job $[s_j, e_j]$. 
\end{mylist}
Each pairwise-disjoint subset is of type (A) or type (B),
so $V(j)$
is the maximum value of any subset of type (A),
or the maximum value of any subset of type (B),
whichever is larger.

As before, the subsets of type (A) are just the pairwise-disjoint
subsets from the first $j-1$ jobs,
so the maximum value of any subset of type (A) is $V(j-1)$.

And, also as before,
each subset of type (B) consists of $[s_j, e_j]$,
together with any pairwise-disjoint subset from
the first $i'-1$ jobs, where $i'$
is the index of the first job that overlaps job $j$.
That is, where $i' = \min\{ i: e_i \ge s_j\}$.
It follows that the maximum value of any subset of type (B) is $V(i'-1) + v_j$.
This gives the recurrence relation.

Correctness of the algorithm (iterative, or recursive and memoized)
% [review 2026-09-27] fix: "the lemma" referred to no stated lemma; correctness follows from the recurrence relation derived above
follows from the recurrence relation by induction on $j$
(that the algorithm does in fact set $W[j]$ to $V(j)$ as previously defined).

\paragraph{run time.}
Here is the iterative algorithm.

\begin{myframed}
  \begin{steps}
  \item[]
    \textsf{maxSchedule}$(C=\{[s_1, e_1], [s_2, e_2], \ldots, [s_n, e_n]\}, v)$:
    \comment{--- iterative version}
    \step for $j=0,1,2,\ldots, n$:
    \step~~~set~\(W[j] = \begin{cases}
      0 & \text{ if } j=0 \\
      \max( W[j-1], W[i'-1] + v_j), \text{ where }  i' = \min\{ i : e_i \ge s_j \}
      & \text{ if } j>0.
    \end{cases}\)
    \step return $W[n]$
  \end{steps}
\end{myframed}

In each iteration, the index $i'$ can be found using binary search
on the jobs (which are ordered by end times) in $O(\log n)$ time.
Then the $n$ iterations take total time $\Theta(n\log n)$.
This proves the following theorem:
\begin{theorem}
  % [review 2026-09-27] fix: theorem named the problem "Weighted Interval Scheduling"; the section's problem is "Weighted Activity Selection"
  There is an $O(n\log n)$-time algorithm for \prob{Weighted Activity Selection}.
\end{theorem}

\paragraph{reduction to Longest Path.}
Consider the DAG that has a vertex for each subproblem $V(j)$,
with edges for each dependency (of weight zero or $v_j$) in the recurrence relation.
The problem is equivalent to finding a max-weight path in this graph from the vertex for $V(0)$
to the vertex $V(n)$.

%%% Local Variables:
%%% mode: latex
%%% TeX-master: "dynamic_programming"
%%% End:
